% RH paper — § RH as an AOR instance
% Drafted body recasting the conditional landing chain as an
% AOR-instance membership statement over the RH carrier.
%
% statements-of-record rows resolved here:
%   def:rh:aor-sel-instance,
%   thm:rh:aor-mechanical-records,
%   thm:rh:aor-recognition-discharge,
%   thm:rh:aor-gamma-bridge-discharge,
%   thm:rh:aor-auditL-discharge,
%   thm:rh:aor-dc-master-import-discharge,
%   thm:rh:aor-real-coordinate-discharge,
%   thm:rh:aor-translation-interface-discharge,
%   thm:rh:aor-instance.

\section{RH as an AOR instance}
\label{sec:rh-aor-instance}

The conditional landing chain of \Cref{sec:landing_chain} may also be
read in the AOR vocabulary of \citet{TsiokosAOR2026}.  The reading is
not a second proof and it does not remove the recognition source.  It
repackages the same closure stack as a membership statement for the
formed carrier \(\Sel\).  In this register, the sentence ``the RH
closure holds on \(\Sel\)'' means that the typed zero ledger, the
functional-equation involution, the trace cone, the recognition source,
the duality-confinement bridge, and translation theorem T form an
audit-closed carrier for the declared RH scope.

The resulting membership judgment has the form
\[
  \Sel \in \mathrm{RefStableAOR}^{\mathrm{ref}}_{\mathfrak S_{\mathrm{RH}}}.
\]
The scope \(\mathfrak S_{\mathrm{RH}}\) is the saturated completed
Selberg trace shell over the typed zero ledger of \(\Lzeta\), restricted
to the typed real-coordinate identification disclosed in
\Cref{sec:involution_and_ledger} and in the formalization appendix.
The record-level reading is that the typed zero ledger of \(\Sel\) has
no nontrivial zeros off \(\Fix(\JL)\).  This is the content of
\Cref{thm:rh:conditional}, read as an AOR-membership witness rather
than as a closure-eliminating derivation.

\subsection{AOR-membership reading of the Selberg carrier}

The AOR meta-theory of \citet{TsiokosAOR2026} treats a carrier as
accepted when its observations, routes, sources, bridges, residuals,
nonclaims, and meta-audits are typed and statused.  Applied here, the
carrier is not the whole classical analytic-number-theoretic universe.
It is the formed shell \(\Sel\) of \Cref{def:rh:sat-sel-shell},
together with its declared zero ledger, trace instrument, involution,
visibility data, and admissibility field.

The membership reading is therefore local to the RH shell.  The
recognition source \(\GamSDTCSelberg\) supplies the completed
domination records
\[
  \AZ\preceq B_n,
  \qquad
  \trace B_n\to0,
\]
as content of the self-dual Selberg trace layer.  The
duality-confinement theorem of \citet{TsiokosSDTC2026} consumes those
records and yields \(\AZ=0\).  Translation theorem T then identifies
\(\AZ=0\) with the critical-line statement on \(\Znt\).  AOR
membership says that this whole typed stack is audit-closed for
\(\mathfrak S_{\mathrm{RH}}\).  It does not assert that the
recognition source has been derived from lower primitives.

\subsection{The Selberg trace shell as an AOR-instance object}

\begin{definition}[The AOR instance object attached to \(\Sel\)]
\label{def:rh:aor-sel-instance}
Let \(\mathfrak S_{\mathrm{RH}}\) be the scope consisting of the
saturated completed Selberg trace shell \(\Sel\), the completed zeta
function \(\Lzeta\), the typed nontrivial-zero ledger \(\Znt\), the
functional-equation involution \(\JL\), the anti-invariant readout
\(\psimRH\), the anti-invariant zero ledger \(\AZ\), the trace
instrument \(\Itr\), the audit field \(\mathrm{Audit}_{L}\), the
recognition source \(\GamSDTCSelberg\), the duality-confinement
master theorem, and translation theorem T.  The AOR instance object
\[
  C_{\mathrm{RH}}(\Sel,\GamSDTCSelberg)
\]
has the following typed fields.
\begin{enumerate}[label=\textup{(\arabic*)}]
  \item The carrier is the zero side of \(\Sel\), with current and
        predictive quotient data inherited from the fields \(Q_L\),
        \(M_L\), and \(\pi_L:M_L\to Q_L\) of
        \Cref{def:rh:sat-sel-shell}.
  \item The observations are the zero ledger \(\Znt\), the
        anti-invariant displacement \(\psimRH\), the zero-ledger
        readout \(\AZ\), and the visibility data \(\mathrm{Vis}_{L}\).
  \item The routes are the functional-equation involution \(\JL\), the
        comparison map \(\pi_L\), the anti-invariant projection
        implicit in the \texttt{AntiInvariantZeroLedger} interface
        reading of \(\psimRH\), and the typed bridge from the
        abstract DC apparatus to the shell.
  \item The source records are the admissibility witness
        \(\mathrm{Audit}_{L}\), the Selberg-class recognition source
        \(\GamSDTCSelberg\), and the completed domination sequence
        \(B_n\).
  \item The interface records are the bridge propositions
        \(\texttt{same\_readout}\), \(\texttt{visible\_zero\_of\_ae}\),
        and \(\texttt{mu\_zero\_of\_ae}\) carried by the
        \(\texttt{GammaSdtcSelberg}\) structure.
  \item The constraints are positivity in the typed cone, the order
        \(\preceq\), trace monotonicity, the DC-ledger domination
        carried by \(\GamSDTCSelberg\), and the asymptotic trace
        condition.
  \item The residual statuses are the discharge records listed in
        \Cref{thm:rh:aor-recognition-discharge,%
        thm:rh:aor-gamma-bridge-discharge,%
        thm:rh:aor-auditL-discharge,%
        thm:rh:aor-dc-master-import-discharge,%
        thm:rh:aor-real-coordinate-discharge,%
        thm:rh:aor-translation-interface-discharge}.
  \item The nonclaim register is the register of
        \Cref{sec:scope_and_nonclaims}, including the conditional
        status of the recognition source, the Foundations-II
        admissibility hypothesis, the typed real-coordinate
        disclosure, and the absence of any GRH or standard-ZFC
        unconditional claim.
\end{enumerate}
\end{definition}

This definition does not add analytic-number-theoretic data.  It
places the paper's existing typed records into the AOR field order.
The carrier is the same shell used in \Cref{thm:rh:conditional}.

\subsection{Mechanical records as membership components}

The definition-side records of the RH paper are mechanical in the AOR
sense.  They introduce typed objects and target records; they do not
supply the recognition content.

\begin{theorem}[Mechanical AOR records for \(\Sel\)]
\label{thm:rh:aor-mechanical-records}
For the scope \(\mathfrak S_{\mathrm{RH}}\), the following records are
AOR-mechanical:
\[
\begin{gathered}
  \text{\Cref{def:rh:fe-involution}, }
  \text{\Cref{def:rh:psi-minus-rh}, }
  \text{\Cref{def:rh:nontrivial-zero-ledger},} \\
  \text{\Cref{def:rh:anti-invariant-zero-ledger}, }
  \text{\Cref{def:rh:sat-sel-shell}, }
  \text{\Cref{thm:rh:translation-T}.}
\end{gathered}
\]
The five definition-side rows discharge by construction.  Translation
theorem T discharges as typed-interface composition over the fields of
the anti-invariant zero ledger.
\end{theorem}

\begin{proof}
The functional-equation involution, the anti-invariant readout, the
nontrivial-zero ledger, the anti-invariant zero ledger, and the
saturated trace shell are introduced as typed records.  They therefore
supply target, role, and presentation data by construction.  Theorem T
then works entirely inside those records: its forward direction reads
\(\AZ=0\) and returns the typed critical-line statement on \(\Znt\);
its reverse direction reads the critical-line statement and returns
\(\AZ=0\).  No recognition-source content is used in the proof of
Theorem T.  Thus the mechanical part of the AOR instance is the
definition-side carrier plus the typed-interface translation theorem.
\end{proof}

\subsection{Substantive and typed-interface discharges}

The remaining records are the records that connect the abstract
duality-confinement apparatus to the Selberg shell and to the RH
readout.  Each receives an AOR residual type and status.

\begin{theorem}[Recognition-source discharge]
\label{thm:rh:aor-recognition-discharge}
The record \(\GamSDTCSelberg\) has primary residual type
\(\Delta_{\mathrm{source}}\).  Its forced secondary residuals are
\(\Delta_{\mathrm{target}}\), \(\Delta_{\mathrm{role}}\), and the
asymptotic trace form of \(\Delta_{\mathrm{limit}}\).  It is
discharged by a bridged recognition-source status: the source supplies
typed completed-domination records \(B_n\) with
\[
  \AZ\preceq B_n,
  \qquad
  \trace B_n\to0.
\]
\end{theorem}

\begin{proof}
The recognition source is the provenance record for the domination
sequence.  Without it, the proof has no source for the inequalities
\(\AZ\preceq B_n\) or for the asymptotic trace ledger.  The target is
\(\AZ\), the role is ``domination source for the self-dual trace
closure,'' and the limit component is the typed convergence
\(\trace B_n\to0\).  These fields are carried by
\(\GamSDTCSelberg\).  The discharge is bridged rather than zero because
the source is supplied as formed-layer recognition content, not
derived from the definition-side fields of \(\Sel\).
\end{proof}

\begin{theorem}[Bridge-field discharge]
\label{thm:rh:aor-gamma-bridge-discharge}
The bridge propositions \(\texttt{same\_readout}\),
\(\texttt{visible\_zero\_of\_ae}\), and
\(\texttt{mu\_zero\_of\_ae}\) carried by
\(\texttt{GammaSdtcSelberg}\) discharge the typed-interface residuals
between the abstract DC apparatus and the shell \(\Sel\).  The primary
residual type is \(\Delta_{\mathrm{transport}}\).  The forced
secondary types are \(\Delta_{\mathrm{role}}\) and
\(\Delta_{\mathrm{target}}\).  The status is bridged, with the
same-readout component discharged at zero.
\end{theorem}

\begin{proof}
The abstract DC apparatus speaks about an involutive object ledger,
its anti-invariant readout, its anti-invariant ledger, and a
measure-theoretic fixed-locus consequence.  The RH shell speaks about
\(\JL\), \(\psimRH\), \(\AZ\), and the typed zero ledger.  The bridge
fields assert that the readouts are the same where the theorem needs
them, that the visible-zero side refines the shell support, and that
the abstract conclusion ``\(\psim\) vanishes \(\mu\)-almost
everywhere'' yields the concrete shell equality \(\AZ=0\).  Thus the
transport from abstract DC output to RH shell output is bridged by the
typed fields, and the role and target mismatches are statused by those
same fields.
\end{proof}

\begin{theorem}[Admissibility-audit discharge]
\label{thm:rh:aor-auditL-discharge}
The \(\mathrm{Audit}_{L}\) field of \(\Sel\) has primary residual type
\(\Delta_{\mathrm{meta}}\).  It is discharged by a bridged
construction-time admissibility witness for the seven Foundations-II
formation schemas.  The discharge does not derive admissibility inside
this paper.
\end{theorem}

\begin{proof}
The shell definition records \(\mathrm{Audit}_{L}\) as the
admissibility field for the saturated completed Selberg trace closure.
AOR membership requires that the carrier's closure apparatus be
meta-audited.  In the present paper that meta-audit is supplied as a
construction-time hypothesis rather than unfolded into derived
predicates.  The correct status is therefore bridged:
\(\mathrm{Audit}_{L}\) carries the admissibility status, while the
paper registers a nonclaim of a separate derivation of that status.
\end{proof}

\begin{theorem}[Cross-paper master-theorem discharge]
\label{thm:rh:aor-dc-master-import-discharge}
The use of the duality-confinement master theorem of
\citet{TsiokosSDTC2026}, realized in Lean by the
\(\texttt{dcMasterApplied}\) bridge to the
\(\texttt{masterTheorem}\) of that source, is an
\(\texttt{approved\_other}\) discharge.  Its primary residual type is
\(\Delta_{\mathrm{transport}}\) from the domination records to the
equality \(\AZ=0\); its secondary types are
\(\Delta_{\mathrm{source}}\), \(\Delta_{\mathrm{role}}\), and
\(\Delta_{\mathrm{target}}\).
\end{theorem}

\begin{proof}
The master theorem is not reproved in this paper.  It is imported as
the abstract typed positive-cone theorem of \citet{TsiokosSDTC2026}.
The imported theorem consumes positive domination records and a
vanishing trace budget and returns the zero anti-invariant ledger.
The source record is the citation to that theorem, the role record is
``duality-confinement membrane theorem,'' and the target record is the
shell equality \(\AZ=0\) after applying the bridge fields of
\Cref{thm:rh:aor-gamma-bridge-discharge}.  Hence the use is
\(\texttt{approved\_other}\) with typed bridge composition, not an
untracked import.
\end{proof}

\begin{theorem}[Typed real-coordinate discharge]
\label{thm:rh:aor-real-coordinate-discharge}
The typed real-coordinate carrier used by the Lean formalization has
primary residual type \(\Delta_{\mathrm{presentation}}\).  The
identification of that carrier with the classical real-part function
\(\Real(\rho)\) is carried by construction-parameter assignment and is
statused as \(\texttt{outside\_scope}/\texttt{nonclaim}\) for a
substantive analytic derivation inside the mathlib-free formalization.
\end{theorem}

\begin{proof}
The paper statement writes \(\Real(\rho)-1/2\).  The formalized
carrier uses a typed \(\texttt{RealCoordinate}\) field supplying the
operations and positivity principles needed by Theorem T.  The
parameter assignment says that, when the shell is instantiated, this
field is the classical real part of the zero.  The present paper does
not derive that analytic identification in Lean.  Therefore the
presentation residual is not ignored: it is closed by the explicit
construction-parameter record together with the nonclaim of an
internal derivation of the classical identification.
\end{proof}

\begin{theorem}[Translation-interface discharge]
\label{thm:rh:aor-translation-interface-discharge}
Translation theorem T has primary residual type
\(\Delta_{\mathrm{target}}\) and secondary type
\(\Delta_{\mathrm{transport}}\).  Both are discharged by construction
over the fields of the anti-invariant zero ledger.  The theorem
transports the typed equality \(\AZ=0\) to the typed critical-line
readout on \(\Znt\), and conversely.
\end{theorem}

\begin{proof}
The proof of Theorem T is a positive-sum argument over the declared
ledger.  The target \(\AZ\) is a finite typed sum of nonnegative
multiplicity-weighted square displacements, and the zero target is
equivalent to the vanishing of every displacement.  Thus the
target-readout and transport obligations are not external analytic
claims; they are the internal typed-interface content of
\Cref{thm:rh:translation-T-forward,thm:rh:translation-T-reverse}.
\end{proof}

\subsection{AOR-instance theorem}

\begin{theorem}[RH AOR-instance theorem]
\label{thm:rh:aor-instance}
Under the discharges recorded in
\Cref{thm:rh:aor-mechanical-records,%
thm:rh:aor-recognition-discharge,%
thm:rh:aor-gamma-bridge-discharge,%
thm:rh:aor-auditL-discharge,%
thm:rh:aor-dc-master-import-discharge,%
thm:rh:aor-real-coordinate-discharge,%
thm:rh:aor-translation-interface-discharge}, the saturated completed
Selberg trace shell is a refinement-stable AOR instance on the RH
scope:
\[
  \Sel\in
  \mathrm{RefStableAOR}^{\mathrm{ref}}_{\mathfrak S_{\mathrm{RH}}}.
\]
Consequently, at the record level of this AOR instance, the typed zero
ledger of \(\Sel\) has no nontrivial zeros off the fixed locus
\(\Fix(\JL)\).
\end{theorem}

\begin{proof}
The definition-side rows provide the carrier, target, role, and
presentation records.  The recognition source provides the
domination-record source.  The \(\texttt{GammaSdtcSelberg}\) bridge
fields carry the abstract DC conclusion to the shell equality
\(\AZ=0\).  The \(\mathrm{Audit}_{L}\) field supplies the admissibility
status of the shell.  The master theorem of \citet{TsiokosSDTC2026}
is an accepted typed-cone import.  The real-coordinate presentation
residual and the admissibility-as-hypothesis residual are both
explicitly statused rather than suppressed.  Translation theorem T
then converts \(\AZ=0\) into the statement
\[
  \forall\,\rho\in\Znt,\quad \Real(\rho)=\frac12.
\]
Thus every residual native to \(\mathfrak S_{\mathrm{RH}}\) is either
zero, bridged, approved, or registered as outside the stronger
formalization scope.  This is precisely the AOR-membership reading of
\Cref{thm:rh:conditional}.
\end{proof}

Lean records this AOR-instance theorem as
\nolinkurl{RH.AORInstance.aorInstance}, returning a value of the local
\nolinkurl{RH.AORPrimitives.RefStableAOR} predicate applied to a
\nolinkurl{RH.AORInstance.assembledCarrier} value built by
concatenating the seven discharge lists.  The proof is by case
analysis on the secondary residual type, exhibiting an explicit
witness atom from the assembled discharge list for each of the seven
local residual types.  The axiom closure is the singleton
\(\{\,\texttt{propext}\,\}\), within the project trust base.  The
Lean theorem returns only the record-level \(\mathrm{RefStableAOR}\)
predicate; the critical-line consequence stated above remains the
conclusion of \Cref{thm:rh:conditional} obtained through the same
closure stack.  The formalization appendix records the
declaration-level correspondence.

\subsection{Partial-status disclosure}

\begin{remark}[Formalization status of the AOR-instance wrapper]
\label{rem:rh:aor-partial-status}
The present AOR-instance theorem is a paper-level classification of
the existing closure records.  The Lean realization carries the
bridge propositions as \(\mathrm{Prop}\)-valued typed hypotheses
inherited from the \(\texttt{GammaSdtcSelberg}\) carrier.  The
AOR-instance proof therefore projects the discharges through the
carrier-attached wrapper: it checks that the carrier has the required
propositions, but it does not inspect data-bearing \(R\)-record values
for those bridges.  The seven per-row discharges of
\Cref{thm:rh:aor-mechanical-records,%
thm:rh:aor-recognition-discharge,%
thm:rh:aor-gamma-bridge-discharge,%
thm:rh:aor-auditL-discharge,%
thm:rh:aor-dc-master-import-discharge,%
thm:rh:aor-real-coordinate-discharge,%
thm:rh:aor-translation-interface-discharge} are realized as
literal-data definitions returning fixed-shape lists of
discharge atoms; substantive Lean composition lives only in
\Cref{thm:rh:aor-instance}.  Promotion of the recognition source,
bridge fields, admissibility witness, and real-coordinate assignment
to data-bearing AOR records is deferred.  This disclosure is part of
the AOR status and is not a weakening of \Cref{thm:rh:conditional}.
\end{remark}
